% titel: Blatt 5, A7: Höhere 1. Ionisierungsenergie

\subsection*{Frage}
Bei welchem Atom je Paar ist die höhere 1.\ Ionisierungsenergie zu erwarten und
warum?

\subsection*{Lösung}
\begin{itemize}
  \item \textbf{Se, Kr:} Kr höher. Gleiche Periode (4), Kr rechts von Se $\to$
        Edelgaskonfiguration ($[Ar]3d^{10}4s^2 4p^6$), grössere $Z_{\mathrm{eff}}$.
  \item \textbf{Kr, Xe:} Kr höher. Gleiche Gruppe (18), Kr einen kleineren
        Atomradius $\to$ Valenzelektronen fester gebunden.
  \item \textbf{O, P:} O höher. P hat die halb besetzte Konfiguration $3p^3$
        (stabil, erniedrigt die IE); O liegt in Periode~2 weiter rechts und hat
        die grössere $Z_{\mathrm{eff}}$ $\to$ O gewinnt.
  \item \textbf{Rb, Sr:} Sr höher. Gleiche Periode (5), Sr rechts von Rb $\to$
        Edelgaskonfiguration des Rumpfes erreicht, höhere $Z_{\mathrm{eff}}$.
  \item \textbf{Li, Cs:} Li höher. Gleiche Gruppe (1), Li oben mit kleinstem
        Radius $\to$ stärkste Anziehung.
\end{itemize}

\begin{ergebnis}
höhere 1.\ IE: Kr (über Se und Xe), O (über P), Sr (über Rb), Li (über Cs).
\end{ergebnis}

\begin{kontrolle}
Alle Aussagen replizieren den bekannten Verlauf: Edelgase haben innerhalb ihrer
Periode die höchste IE (Kr rechts von Se), und in ihrer Gruppe ist das oberste
Element am schwersten zu ionisieren (Kr über Xe, Li über Cs). $\mathrm{O}>\mathrm{P}$
liegt an der Lage O weiter rechts trotz $3p^3$-Stabilisierung von P --- die
grössere $Z_{\mathrm{eff}}$ gewinnt. Referenzreihenfolge: $\mathrm{Rb<Sr<Kr}$,
$\mathrm{Li<Na<\dotsb}$.
\end{kontrolle}

\begin{merksatz}
Höhere 1.\ IE: rechts in der Periode (Edelgas), oben in der Gruppe (kleines
Atom). Der Radius-Trend zeigt genau umgekehrt: kleiner Radius $\leftrightarrow$
höhere IE.
\end{merksatz}
